\documentclass[11pt,a4paper]{article}
\usepackage[T1]{fontenc}
\usepackage[utf8]{inputenc}
\usepackage[margin=25mm]{geometry}
\usepackage{lmodern}
\usepackage{microtype}
\usepackage{booktabs,tabularx,longtable,array}
\usepackage{xurl}
\usepackage[colorlinks=true,linkcolor=blue,urlcolor=blue,citecolor=blue]{hyperref}
\usepackage[backend=biber,style=numeric,sorting=none,maxnames=3,minnames=1,
  doi=true,url=true,isbn=true]{biblatex}
\addbibresource{thesis_reference_review.bib}
\setlength{\parindent}{0pt}
\setlength{\parskip}{5pt}
\setcounter{secnumdepth}{1}
\setcounter{tocdepth}{2}
\newcommand{\referenceentry}[3]{%
  \subsection{#1: #2}\label{ref:#3}%
  \fullcite{#3}\par
  \textbf{Citation key:} \texttt{#3}.\par}
\newcommand{\field}[2]{\textbf{#1:} #2\par}

\title{LM-IDNet: Verified References for Thesis Writing\\
  \large An annotated shortlist with claim boundaries and reading pointers}
\author{Project reference review}
\date{Verified 16 October 2026}

\begin{document}
\maketitle

This guide contains \textbf{30 scholarly works}: 28 journal or conference
publications, one technical report, and one book. Two additional dataset records
are listed separately. The selection follows the actual project: four-category
traffic windows, a Dirichlet--multinomial (DM) device model, Minka estimation,
LM likelihood evaluation, normal-data calibration, guarded adaptation,
chronological evaluation, and a resource-limited Docker simulation.

The annotations are reading and attribution notes, not ready-made thesis
chapters. Each explains what a source can support, where it belongs, and what
it cannot establish. \textbf{A reference supports an idea or precedent; it does
not validate this project's numerical results or establish its novelty.}

\tableofcontents
\clearpage

\section{How to use this shortlist}

\subsection{What was verified}

Titles, authors, publication venues, years, page ranges or article identifiers,
and DOI/official URLs were checked against publisher pages, author or university
records, proceedings, and DOI registrant metadata. Existing PDFs in
\path{resouces/} were examined rather than trusting their filenames. Final
publication metadata takes precedence over provisional manuscript headers.

\textbf{Full text} means the relevant passage was inspected in a publisher
copy, an author manuscript, a local supplied PDF, or an authorised HTML reprint.
It does not mean every proof and experiment was audited. There are 27 such
entries, with the particular version identified where necessary.
\textbf{Limited access} applies to R01, R18, and R25: their existence and
bibliographic identity were verified, but detailed passage-level checking was
limited to the publisher record, abstract, or chapter summary. No unverified
equation or narrow page claim is attributed to them here.

The suggested use of each reference is an interpretation of its relevance to
LM-IDNet. Page pointers distinguish printed journal pages from PDF page numbers.
An author manuscript can have different pagination and content from the final
publication. Search-result dates were not treated as publication dates.

\subsection{Corrections to existing project references}

\begin{itemize}
\item The IoT forecasting study called a ``reference manuscript'' in
  \path{docs/full_project_report/bibliography.tex} is published at
  \textbf{ICNC 2025, pp.~28--32}
  (\href{https://doi.org/10.1109/ICNC64010.2025.10993688}{published record}, R06).
  Use that publication when attributing the four-category,
  ten-minute Camera~5 precedent.
\item \path{resouces/ALM2_new.pdf} is a 15-page preprint. The matching title and
  authors have a published \textbf{ACDSA 2025, pp.~1--6} record (R07).
  The final article and long preprint must not be assumed identical.
\item The LM article belongs to the \textbf{2023} volume of
  \emph{Computational Statistics}, although it appeared online in 2022 (R04).
  The TPAMI article belongs to the \textbf{2025} issue, although its DOI and
  online publication date contain 2024 (R05).
\item The UNSW attack-trace page explicitly requests citation of
  \textbf{Hamza et al., SOSR 2019} (R13). This is a different source from
  UNSW-NB15, N-BaIoT, and the UNSW benign device-classification traces.
\item Minka's report is \textbf{2000, revised through 2012} (R02), not a journal
  article. Its DM count estimator is Eq.~(55), not the Dirichlet-proportions
  update from its first section.
\end{itemize}

These corrections are recorded here; existing thesis chapters and older reports
have not been rewritten.

\subsection{A chapter and claim map}

\begin{longtable}{@{}>{\raggedright\arraybackslash}p{0.43\textwidth}
  >{\raggedright\arraybackslash}p{0.50\textwidth}@{}}
\toprule
\textbf{Topic or claim to support} & \textbf{Start with these entries} \\
\midrule
\endhead
Introduction: IoT compromise and motivation &
R10 (Mirai), R11 (device traffic), R28 (edge context). \\
Background: anomaly detection and normal-reference assumptions &
R08 (taxonomy), R09 (operational limitations), R20 (training contamination). \\
Background: categorical counts and extra-multinomial variation &
R01 (history), R02--R03 (accessible technical foundations), R05. \\
Methodology: estimator versus numerical likelihood backend &
R02 (update), R03--R05 (numerical methods), R07 and R26--R27 (implementation context). \\
Methodology: inherited IoT representation and forecasting precedent &
R06; use D1 separately for the original D-Link data. \\
Related work: device profiles, shallow and neural detectors &
R11--R17. R13 is also the attack-dataset paper. \\
Adaptation: changing normal behaviour and contamination risks &
R12, R18--R20. These motivate safeguards, not their exact constants. \\
Experimental methodology: temporal splits, leakage, and development decisions &
R21--R22; document the actual capture selection locally. \\
Experimental methodology: imbalance and time-dependent anomalies &
R23 (precision/recall), R24 (range metrics and timing). \\
Experimental methodology: calibration quantiles and reproducibility &
R30 (quantile convention), R29 (workflow provenance). \\
Statistical comparisons: KL divergence &
R25; R05's Table III also documents the inherited distance measures. \\
Datasets, devices, annotations &
R13 plus D2 for UNSW; R06 plus D1 for D-Link. \\
Discussion: edge feasibility and limits &
R15, R17, R28 for context; this project's measured hardware evidence separately. \\
\bottomrule
\end{longtable}

For a first reading pass, prioritise R02, R04--R06, R08--R09, R12--R15,
R19, and R21--R23. Read the remaining items where the corresponding chapter
needs them; a shortlist is not a requirement to cite everything.

\section{Statistical foundations and the inherited LM work}

\referenceentry{R01}{Historical compound-multinomial foundation}{mosimann1962}
\field{Use and placement}{Background: historical attribution for compounding a
multinomial with a distribution over proportions. It prevents presenting the
DM family as a new model introduced by this thesis.}
\field{Verification and reading}{\textbf{Limited access.} Oxford's publisher
record confirms the title, author, 1962 volume, and pp.~65--82. The article is
PDF-only behind access controls; no specific derivation page was verified.
Use R02--R03 for the technical passages available during this review.}
\field{Boundary}{A foundational probability paper, not evidence that a
particular IoT trace is overdispersed or that DM is the best detector.}
\field{Project connection}{\path{docs/full_project_report/sections/scientific_basis.tex}:
the DM generative model and concentration interpretation.}

\referenceentry{R02}{The parameter estimator actually used}{minka2000}
\field{Use and placement}{Background and methodology: maximum-likelihood
estimation of DM parameters using fixed-point iteration. Attribute the alpha
update to Minka, including when the likelihood backend is LM.}
\field{Verification and reading}{\textbf{Full text.} Author's publication page
and local \path{resouces/minka-dirichlet.pdf}. Section~3, PDF p.~5,
Eqs.~(48)--(55); Eq.~(55) is the project's digamma-based update. Appendix~B,
PDF pp.~12--13, gives its derivation.}
\field{Boundary}{The report also models ordered categorical outcomes. Its
likelihood expressions do not justify omitting the multinomial coefficient
from a probability assigned to a count vector. A parameter-independent term
can be dropped for fitting and still be necessary for cross-window scoring.}
\field{Project connection}{\path{src/lm_idnet/algorithms/dirichlet.py} and
\path{docs/likelihood-semantics.md}.}

\referenceentry{R03}{The Yu--Shaw numerical predecessor}{yu2014}
\field{Use and placement}{Background: the numerical problem near the
multinomial limit and the Yu--Shaw approach to evaluating the DM likelihood.
This is a mathematical predecessor even though its application is genomics.}
\field{Verification and reading}{\textbf{Full text.} Local
\path{resouces/bioinformatics_30_11_1547.pdf}, especially the introduction
on printed p.~1547 and the methods that follow. PubMed corroborates the
authors, journal, pp.~1547--1554, and DOI.}
\field{Boundary}{Its speed and accuracy results concern its own numerical
tests. They do not describe this project's four-category detector or establish
an LM-versus-SciPy speed advantage.}
\field{Project connection}{The Yu--Shaw comparison in
\path{thesis/chapters/02_background_and_related_work.tex} and the numerical
discussion in \path{docs/lm_paper_project_comparison_report.md}.}

\referenceentry{R04}{The Languasco--Migliardi algorithm}{languasco2023}
\field{Use and placement}{Background and methodology: the LM evaluation of
nearby log-Gamma differences, error control, and computational complexity.
This is the principal attribution for the native numerical kernel.}
\field{Verification and reading}{\textbf{Full text.} Local
\path{resouces/180_2022_1311_OnlinePDF.pdf} is a revised proof, not the
paginated final issue. Read its introduction, Section~3 (Theorem~1), and
Sections~4--5 (complexity and experiments). Padua's record and author
publication list confirm the final 2023 citation, pp.~1995--2013.}
\field{Boundary}{LM evaluates likelihood terms; it does not supply the
Minka update or the project's anomaly decision rule. Numerical improvements
must be separated from detection accuracy.}
\field{Project connection}{\path{src/lm_idnet/algorithms/lm_native.py},
\path{src/lm_idnet/algorithms/native/}, and the native-backend chapter.}

\referenceentry{R05}{DM estimation and cross-domain numerical evidence}{alhajbaddar2025analysis}
\field{Use and placement}{Background and methodology: DM fitting for
overdispersed counts, fixed-point estimation, and empirical comparison of
LM and Yu--Shaw. Also a direct precedent for the project's distance measures.}
\field{Verification and reading}{\textbf{Full text.} Local TPAMI PDF:
Section~II and Algorithm~1 on p.~1182; experiments in Section~III;
Table~III on p.~1186 lists TVD, KL divergence, Euclidean distance, and MSE.
The issue header and publication record agree on 47(2), 1181--1189, 2025.}
\field{Boundary}{Polls, images, and IoT examples support a numerical modelling
study. Its reported speedups and goodness-of-fit results are not LM-IDNet
attack-detection scores or a universal comparison with SciPy.}
\field{Project connection}{\path{docs/full_project_report/sections/numerical_experiments.tex}
and \path{src/lm_idnet/evaluation/metrics.py}.}

\referenceentry{R06}{The direct four-category IoT precursor}{alhajbaddar2025iot}
\field{Use and placement}{Background and methodology: attribution for D-Link
Camera~5 counts, TCP/UDP/SSDP/ARP categories, ten-minute aggregation,
DM fitting, and traffic-profile forecasting. This is the closest source for
the inherited representation.}
\field{Verification and reading}{\textbf{Full text.} Local
\path{resouces/ALM_IoT-CR.pdf}, Section~III beginning on PDF p.~2;
modelling and forecast tables on later pages. Padua's publication record,
the author's publication list, and DOI metadata confirm ICNC 2025,
pp.~28--32. The local copy supplies the reading pagination.}
\field{Boundary}{This study forecasts traffic; it does not establish the
project's calibrated attack detection, guarded updates, or precision/recall.
Ten minutes is a precedent, not proof of an optimal window duration.}
\field{Project connection}{\path{docs/lm_paper_project_comparison_report.md}
and \path{docs/full_project_report/sections/window_size_experiments.tex}.}

\referenceentry{R07}{The later C-wrapped LM implementation}{alhajbaddar2025implementation}
\field{Use and placement}{Methodology and related numerical work: a
Python-accessible C implementation of LM and comparison with standard
Python/SciPy likelihood evaluation. It helps attribute the supplied code's
implementation lineage.}
\field{Verification and reading}{\textbf{Full text of the supplied preprint.}
\path{resouces/ALM2_new.pdf}: Section~3 discusses the C wrapper and Section~4
the experiments. The 15-page preprint is dated 11 April 2025. Padua's
record, the author list, and DOI metadata identify the later six-page ACDSA
2025 publication. Final-text equivalence was not checked.}
\field{Boundary}{Do not transfer the preprint's page numbers, dataset counts,
or speedups to the final article without comparing versions. LM-IDNet's
backend benchmark needs its own measurements.}
\field{Project connection}{\path{lm-algorithm/}, the native wrapper, and
\path{docs/lm_backend_change_record.tex}.}

\section{Anomaly detection, IoT behaviour, and related systems}

\referenceentry{R08}{Taxonomy and statistical anomaly detection}{chandola2009}
\field{Use and placement}{Introduction and background: define anomalies,
explain statistical low-likelihood detection, and distinguish supervised,
normal-reference, and unlabelled-data settings.}
\field{Verification and reading}{\textbf{Full text.} Author-hosted published
article; Section~2 covers problem characteristics and training labels;
Section~7 covers statistical techniques. The first page confirms article~15,
41(3), 2009, 58 pages.}
\field{Boundary}{In this taxonomy, training on presumed normal data is commonly
called semi-supervised anomaly detection. Define precisely why this project
uses ``unsupervised'': fitting and calibration require no attack labels, but
a trusted normal reference is still assumed. Do not claim that labels never
enter method selection or offline evaluation.}
\field{Project connection}{The README's learning assumptions and
\path{thesis/chapters/02_background_and_related_work.tex}.}

\referenceentry{R09}{Operational limits of learned network detectors}{sommer2010}
\field{Use and placement}{Introduction, related work, and limitations:
motivate the distinction between a laboratory anomaly detector and a useful
operational intrusion detector. Discuss normal-traffic diversity, evaluation,
false alerts, and understanding decisions.}
\field{Verification and reading}{\textbf{Full text.} Published 12-page paper
inspected through a university-hosted
\href{https://gangw.cs.illinois.edu/class/cs562/papers/Closed_World-IDS-sp10.pdf}{copy};
read the abstract, introduction, and discussion of challenges and guidelines.
DOI metadata confirms IEEE S\&P 2010, pp.~305--316.}
\field{Boundary}{A general discussion of network intrusion detection, not a
proof that machine learning cannot work or that a small DM model is superior.}
\field{Project connection}{The trusted-baseline assumption, alert evidence,
and deployment limits in \path{docs/full_project_report/sections/interpretation.tex}.}

\referenceentry{R10}{An observed IoT botnet threat}{antonakakis2017}
\field{Use and placement}{Introduction: motivate IoT security using a measured
case of embedded devices participating in a large DDoS botnet.}
\field{Verification and reading}{\textbf{Full text.} Official USENIX paper and
proceedings citation; introduction on printed p.~1093 and the measurement
analysis that follows. The PDF includes an extra proceedings cover, so PDF
page numbers are offset from article pages.}
\field{Boundary}{Mirai is a historical case study. Do not turn its measurements
into present-day device totals, claim every IoT attack is volumetric, or say
this project was evaluated on Mirai unless the dataset actually supports it.}
\field{Project connection}{The threat motivation; the project's UNSW labels
describe their own attacks and are not Mirai ground truth.}

\referenceentry{R11}{Device-specific network traffic characteristics}{sivanathan2019}
\field{Use and placement}{Background: empirical evidence that consumer IoT
devices can exhibit useful device-specific traffic attributes, including
activity cycles, flow volume, ports, and service interactions.}
\field{Verification and reading}{\textbf{Full text.} Published IEEE copy and
author's publication list; read the traffic-characterisation analysis,
especially the attribute summary on printed p.~1752. Issue pagination is
1745--1759, August 2019; online publication and some author BibTeX use 2018.}
\field{Boundary}{This is primarily device classification using richer features.
It does not prove that four protocol counts identify every attack or that
one device's normal profile transfers to another.}
\field{Project connection}{Per-device fitting and the compact representation
described in \path{docs/full_project_report/sections/scientific_basis.tex}.}

\referenceentry{R12}{Changing device behaviour and one-class models}{sivanathan2020}
\field{Use and placement}{Related work: per-device one-class traffic models
and detection of behavioural change, including legitimate changes associated
with firmware upgrades. This is especially close to the motivation for
adaptation.}
\field{Verification and reading}{\textbf{Full text of the author manuscript.}
Sections~III--IV describe flow attributes and clustering; Section~V.B--C
evaluates behavioural changes and attacks. DOI metadata verifies the final
7(8), 7295--7309, 2020 citation; manuscript pagination differs.}
\field{Boundary}{Its clustering/classification architecture differs from DM
likelihood scoring. A detected change can be benign; its update mechanisms
and performance are not those implemented in LM-IDNet.}
\field{Project connection}{\path{docs/full_project_report/sections/conservative_adaptation.tex}
and the related-systems chapter.}

\referenceentry{R13}{The UNSW attack-dataset paper and volume monitoring}{hamza2019}
\field{Use and placement}{Related work and dataset methodology: provenance
for the UNSW attack traces; monitoring per-device traffic activity to detect
volumetric attacks and ARP spoofing. It is also relevant to the separate total
and ARP count signals.}
\field{Verification and reading}{\textbf{Full text of the author copy.}
Sections~3--4 describe MUD/SDN monitoring and attack/data collection;
Section~5 evaluates the prototype. DOI metadata confirms final pp.~36--48;
the author copy numbers its pages 1--13. The dataset page names this paper
as its requested citation (D2).}
\field{Boundary}{LM-IDNet uses protocol windows, not this paper's MUD flow
rules, SDN pipeline, or classifier. Cite the paper and actual annotations;
do not describe the selected traces as UNSW-NB15.}
\field{Project connection}{\path{docs/unsw-chromecast-dataset.md},
\path{docs/unsw-samsung-camera-dataset.md}, and joint count scoring.}

\referenceentry{R14}{Normal-device learning with deep autoencoders}{meidan2018}
\field{Use and placement}{Related work: N-BaIoT learns benign device
behaviour and detects deviations using deep autoencoders. Compare its feature
representation and learning model with the project's compact count model.}
\field{Verification and reading}{\textbf{Full text of the author manuscript.}
Sections~3--4 describe the detector and evaluation; the feature discussion
describes statistics over multiple temporal windows. Publisher DOI metadata
confirms \emph{IEEE Pervasive Computing} 17(3), pp.~12--22 and the author
order used in the companion BibTeX. The arXiv header/order differs.}
\field{Boundary}{Its nine-device Mirai/BASHLITE study and dataset are distinct
from the UNSW traces used here. Published detection figures cannot be ranked
directly against LM-IDNet's results on different traces and window labels.}
\field{Project connection}{The neural-detector comparison in related work;
not an implemented experimental baseline.}

\referenceentry{R15}{Online and resource-aware network detection}{mirsky2018}
\field{Use and placement}{Related work: Kitsune/KitNET combines an ensemble
of small autoencoders with incremental feature extraction for online network
anomaly detection. It is a concrete precedent for gateway-oriented detection.}
\field{Verification and reading}{\textbf{Full text.} Official NDSS paper;
Section~IV describes the pipeline and complexity, Section~V the experiments,
and Section~VI adversarial issues. The NDSS programme and DOI record confirm
the title, authors, and 2018 venue.}
\field{Boundary}{Kitsune's Raspberry Pi measurements belong to Kitsune.
They do not establish the hardware performance of this project's Docker
simulation. Neither unsupervised online detection nor low-resource deployment
is a new idea introduced by this thesis.}
\field{Project connection}{Bounded streaming and the comparison of count-based
and neural representations.}

\referenceentry{R16}{Device-type profiles and federated self-learning}{nguyen2019}
\field{Use and placement}{Related work: D\"IoT models device-type-specific
communication sequences and combines learning across gateways. It shows
another way to learn normal IoT behaviour without attack-labelled training.}
\field{Verification and reading}{\textbf{Full text of the accepted manuscript.}
The arXiv copy identifies acceptance at ICDCS 2019. Read Section~III
(packet-sequence modelling), Section~IV (federated learning), and
Section~VII.B--D (behaviour evolution and adversarial limitations).
Aalto's record and DOI metadata confirm pp.~756--767.}
\field{Boundary}{LM-IDNet has neither federated training nor a GRU sequence
model. Use D\"IoT as a related system, not as attribution for the specific
DM update gate or evidence of transfer between devices.}
\field{Project connection}{Device-specific baselines and the learning-architecture
comparison in background.}

\referenceentry{R17}{Compact IoT features and gateway placement}{doshi2018}
\field{Use and placement}{Related work and motivation: selecting traffic
features using IoT communication properties and evaluating DDoS detection
with several machine-learning classifiers at a network middlebox.}
\field{Verification and reading}{\textbf{Full text.} IEEE conference-hosted
published PDF: introduction, feature-engineering subsection, and results.
Printed pp.~29--35 and DOI 10.1109/SPW.2018.00013 are visible in the
publisher copy. DLS was the workshop; SPW is the proceedings citation.}
\field{Boundary}{The study uses labelled classifier experiments and features
different from the project's four counts. It motivates low-cost monitoring,
but neither proves an optimal feature set nor demonstrates this detector's
physical gateway performance.}
\field{Project connection}{Representation size and gateway placement; a
literature comparison, not an algorithm included in the current experiments.}

\section{Adaptation, drift, and contaminated learning data}

\referenceentry{R18}{Concept drift terminology and adaptation context}{gama2014}
\field{Use and placement}{Background: established terminology for evolving
data and adaptive learning, and a framework for introducing why a static
normal model may become stale.}
\field{Verification and reading}{\textbf{Limited access.} Publisher abstract
and bibliographic record checked: ACM Computing Surveys 46(4), article~44,
37 pages, 2014. The abstract explicitly describes its main setting as
\emph{online supervised learning}. Obtain the full final text before relying
on a detailed algorithm or page-specific argument.}
\field{Boundary}{Its central setting includes target labels. In LM-IDNet,
describe changing unlabelled traffic as distributional or behavioural drift
unless a stronger concept-drift claim is justified. The survey does not
validate the project's update policy or update constants.}
\field{Project connection}{The rationale for static, threshold-adaptive, and
periodic-refit modes.}

\referenceentry{R19}{Poisoning an online anomaly detector}{kloft2012}
\field{Use and placement}{Background, adaptation methodology, and limitations:
online anomaly models can be manipulated by admitted observations; safeguards
and attacker assumptions need to be explicit.}
\field{Verification and reading}{\textbf{Full text.} JMLR's official article
and 44-page PDF, pp.~3681--3724. The abstract and attack analysis distinguish
full/limited traffic control and external false-positive constraints.}
\field{Boundary}{The formal bounds concern online centroid anomaly detection,
not DM fitting. Cite the poisoning risk as motivation; do not claim those
bounds prove the safety of this project's confidence gate, fixed baseline,
or per-update movement limits. Small accepted shifts may accumulate.}
\field{Project connection}{\path{docs/full_project_report/sections/conservative_adaptation.tex}:
admission gates, bounded buffers, fixed references, and movement checks.}

\referenceentry{R20}{Why unlabelled training data may need sanitisation}{cretu2008}
\field{Use and placement}{Background and methodology: contaminated training
traffic can corrupt anomaly sensors. This supports stating the trusted-normal
assumption and motivates selective admission during adaptation.}
\field{Verification and reading}{\textbf{Full text of the author copy.}
Sections~2--3 describe local sanitisation and its evaluation; Section~3.4
discusses long-lasting training attacks; Section~4 introduces collaborative
sanitisation. The paper and DOI record identify IEEE S\&P 2008,
pp.~81--95.}
\field{Boundary}{The paper's micro-model voting and collaborative cleaning
are different mechanisms. LM-IDNet's filtering is an engineering safeguard,
not an implementation of this sanitisation algorithm or a guarantee that
accepted observations are attack-free.}
\field{Project connection}{Presumed-benign fitting/calibration inputs and
the conservative learning buffer.}

\section{Evaluation design and interpretation}

\referenceentry{R21}{Temporal and distributional evaluation bias}{pendlebury2019}
\field{Use and placement}{Experimental methodology: motivate chronological
training/testing and deployment-representative distributions. Explain why
random mixing across time can produce misleading security results.}
\field{Verification and reading}{\textbf{Full text.} Official USENIX paper,
Sections~3.2--3.3 and~4.1. Its
\href{https://www.usenix.org/system/files/sec19_errata_6.pdf}{official erratum}
corrects an extra normalisation factor in Eq.~(4) on p.~735;
the authors state their released implementation and results are unaffected.}
\field{Boundary}{The experiments concern Android malware classification.
The evaluation principle transfers, but the paper does not verify this
repository's splits. Earlier label exposure still prevents a later partition
from becoming an independently locked test merely because it is chronological.}
\field{Project connection}{\path{docs/temporal-partitions.md} and
\path{docs/full_project_report/sections/datasets.tex}.}

\referenceentry{R22}{Security-specific experimental pitfalls}{arp2022}
\field{Use and placement}{Experimental methodology and limitations: justify
clear baselines, label provenance, leakage controls, appropriate metrics,
and careful separation of laboratory measurements from deployment claims.}
\field{Verification and reading}{\textbf{Full text.} Official USENIX paper;
Figure~1 and Section~2 identify data snooping, biased parameters,
inappropriate measures, base-rate effects, and lab-only evaluation.
The proceedings page confirms pp.~3971--3988.}
\field{Boundary}{This paper supplies methodological guidance. It cannot certify
that LM-IDNet avoids every pitfall. Any development-label use for choosing
window length, count signals, or adaptation policy must be disclosed even
when labels never enter the runtime learning rule.}
\field{Project connection}{\path{thesis/chapters/04_experimental_methodology.tex}
and the project's explicit limitations.}

\referenceentry{R23}{Precision and recall under class imbalance}{saito2015}
\field{Use and placement}{Experimental methodology: explain precision, recall,
false positives, and why accuracy or an ROC plot alone can hide the operational
burden when attack windows are rare.}
\field{Verification and reading}{\textbf{Full text.} PLOS publisher article;
``Theoretical Background,'' Table~1, and the subsection on precision for
imbalanced datasets. The publisher confirms 10(3), e0118432, 2015.}
\field{Boundary}{The motivating datasets are biological, but the binary-metric
argument applies to labelled windows. It does not say ROC analysis is always
invalid. Window precision is not attack-episode recall; class prevalence and
the unit of evaluation must remain explicit.}
\field{Project connection}{Confusion counts, precision/recall/F1, and the
class-imbalanced Chromecast and Samsung development partitions.}

\referenceentry{R24}{Time ranges and event-level detection}{tatbul2018}
\field{Use and placement}{Experimental methodology and limitations: motivate
why detecting a time interval has additional dimensions such as coverage,
fragmentation, and positional preference, beyond pointwise classification.}
\field{Verification and reading}{\textbf{Full text.} Author-hosted paper and
official NeurIPS 2018 record. Section~4 defines range precision/recall;
Section~4.3 discusses customisation. The official record identifies volume~31;
no unverified proceedings page range is inserted in the BibTeX.}
\field{Boundary}{The project's window metrics and episode summary are not
automatically Tatbul's range metrics. Cite it as motivation and define the
implemented metric separately. A label triggered by any attack overlap is
different from weighting a window by its attack duration.}
\field{Project connection}{Window-size comparisons, annotation overlap,
episode interpretation, and response-time reporting.}

\section{Supporting mathematics, software, deployment, and reproducibility}

\referenceentry{R25}{A book source for relative entropy}{cover2006}
\field{Use and placement}{Background or methodology: information-theoretic
context for KL divergence when comparing fitted and observed category
distributions. Cite this book for the definition, and R05--R06 for its use
in the inherited modelling experiments.}
\field{Verification and reading}{\textbf{Limited access.} Wiley's book record
confirms the second edition and ISBN. Its table of contents assigns
Chapter~2, ``Entropy, Relative Entropy, and Mutual Information,''
\textbf{pp.~13--55}; the
\href{https://onlinelibrary.wiley.com/doi/10.1002/047174882X.ch2}{chapter summary}
explicitly lists relative entropy. The chapter's complete text was not inspected.
Use the 2006 print-edition year; Wiley also exposes a 2005 online date.}
\field{Boundary}{No narrower book page or theorem was verified. Do not cite
this chapter for a specific Jensen--Shannon or TVD formula without checking it.
KL between category profiles is not an attack probability.}
\field{Project connection}{\path{src/lm_idnet/evaluation/metrics.py}.}

\referenceentry{R26}{Floating-point cancellation and numerical stability}{goldberg1991}
\field{Use and placement}{Background: explain why subtracting nearby
large floating-point values can expose rounding error, motivating stable
formulations of log-Gamma differences and reference calculations.}
\field{Verification and reading}{\textbf{Full text.} Oracle's authorised,
edited HTML reprint, the ``Cancellation'' subsection under ``Rounding
Error.'' The reprint identifies the original March 1991 ACM article;
its original citation is 23(1), pp.~5--48.}
\field{Boundary}{The HTML reprint's layout/pagination is not the original
journal pagination. This is general numerical background; R03--R04 provide
the DM-specific argument. It does not prove a particular error bound for
the project's compiled kernel.}
\field{Project connection}{Backend accuracy checks and
\path{docs/full_project_report/sections/numerical_experiments.tex}.}

\referenceentry{R27}{Credit for the numerical software dependency}{virtanen2020}
\field{Use and placement}{Methodology and reproducibility: acknowledge SciPy,
including the special functions used by the estimator and the reference
likelihood implementation.}
\field{Verification and reading}{\textbf{Full text.} Nature Methods publisher
article, the subpackages overview and its \texttt{special} description:
Gamma-related and information-theory functions are listed. Publisher and DOI
metadata confirm 17(3), pp.~261--272, 2020. The publisher also links an author
correction; this guide relies on the current article's software description.}
\field{Boundary}{The paper is a software citation, not documentation of the
exact installed version or a guarantee of accuracy in every regime. Record
the project's actual package versions and benchmark environment separately.}
\field{Project connection}{\path{pyproject.toml},
\path{src/lm_idnet/algorithms/dirichlet.py}, and
\path{src/lm_idnet/algorithms/log_likelihood.py}.}

\referenceentry{R28}{Why process data near the network edge}{shi2016}
\field{Use and placement}{Introduction and deployment background: define
edge computing and motivate local processing in terms of latency, bandwidth,
resource constraints, and data handling.}
\field{Verification and reading}{\textbf{Full text.} Author-hosted published
IEEE paper, Section~II, ``What Is Edge Computing,'' and its case studies and
challenges. The first page confirms 3(5), October 2016, p.~637;
the article ends at p.~646 and prints the DOI.}
\field{Boundary}{Architectural motivation is not evidence of LM-IDNet running
on physical IoT hardware, meeting packet-line-rate requirements, or reducing
energy use. A one-CPU, 512-MiB Linux/amd64 Docker simulation supports the
measurements actually collected in that environment.}
\field{Project connection}{\path{docs/edge-deployment.md} and
\path{docs/full_project_report/sections/edge_deployment.tex}.}

\referenceentry{R29}{Reproducible computational experiments}{sandve2013}
\field{Use and placement}{Experimental methodology and reproducibility
appendix: motivate tracking inputs, commands, software versions, random seeds,
intermediate artifacts, and the path from raw data to a result.}
\field{Verification and reading}{\textbf{Full text.} PLOS publisher article;
Rules~1--4 cover workflow, manual steps, software versions, and version
control; Rules~6--7 cover random seeds and intermediate results. The publisher
confirms 9(10), e1003285, 2013.}
\field{Boundary}{This is guidance originating in computational research,
not a claim that fingerprints or containers alone guarantee reproducibility.
The actual artifacts, versions, and commands remain part of the thesis's
own evidence.}
\field{Project connection}{\path{docs/dataset-fingerprints.md},
\path{docs/deterministic-randomness.md}, and saved experiment manifests.}

\referenceentry{R30}{The empirical quantile convention}{hyndman1996}
\field{Use and placement}{Methodology: empirical quantiles have multiple
definitions, so calibration should state its interpolation convention.
This is a reference for the estimator of the quantile, not for selecting
the one-percent operating point.}
\field{Verification and reading}{\textbf{Full text.} Author's five-page scanned
PDF; printed p.~363 was rendered and inspected. Definition~7 uses plotting
positions $(k-1)/(n-1)$. The
\href{https://numpy.org/doc/stable/reference/generated/numpy.quantile.html}{official NumPy documentation}
also explicitly maps \texttt{method="linear"} to Definition~7.
The author's publication page confirms 50(4), pp.~361--365 and lists
\href{https://robjhyndman.com/publications/quantiles/}{errata}, including
corrections to the properties and tail-frequency discussion.}
\field{Boundary}{A one-percent empirical lower quantile is not a guaranteed
one-percent future false-positive rate under dependence or drift. The paper
does not choose the project's quantile, IQR gates, or movement bounds.}
\field{Project connection}{\path{src/lm_idnet/models/calibration_stage.py}:
\texttt{np.quantile(..., method="linear")}, calibration artifacts, and IQR scales.}

\section{Supplementary dataset records}

These two records are \textbf{additional to the 30 scholarly works}. Cite
the dataset and the relevant method paper together when describing provenance.

\referenceentry{D1}{Original D-Link dataset, version 1}{chowdhury2021data}
\field{Use and placement}{Dataset methodology: original source of D-Link
packet traces used in the Camera~5 modelling workload.}
\field{Verification and reading}{Mendeley Data's publisher page confirms
the four contributors, version~1, 14 January 2021 publication date, DOI,
14-device description, and CC~BY~4.0 licence. It describes packet-level and
802.11 MAC-frame traces, collected between September 2020 and January 2021.}
\field{Boundary}{The dataset description is not attack ground truth for the
project's selected Camera~5 captures. The selected dates, MAC filter,
counts, partitions, and fingerprints must be documented from project records.
R06 is the source for the subsequent four-category modelling procedure.}

\referenceentry{D2}{UNSW attack traces and annotation description}{unswattackdata}
\field{Use and placement}{Dataset methodology: publisher's release of benign
and mixed attack/benign daily packet traces and device attack annotations.
The page's ``Cite our data'' block requests R13.}
\field{Verification and reading}{The web index exposed the publisher's original
description, date lists, annotation example, and citation block. Direct HTTP
retrieval returned 404 during this review, so the link is recorded as a
\textbf{source with an access limitation}, not a freshly downloadable release.
R13's accessible author paper independently describes the data collection.}
\field{Boundary}{The publisher's example labels its own flow records;
LM-IDNet's half-open-window overlap rule is a separate processing choice.
The repository's annotation hashes and coverage/exclusion records establish
which bytes and intervals were actually used. The local licence statement
was not independently reconfirmed from a live release page.}

\section{Attribution boundaries to retain while drafting}

\begin{itemize}
\item \textbf{Inherited ideas:} DM modelling and fixed-point estimation
  (R01--R03); LM likelihood evaluation (R04--R05, R07); Camera~5,
  the protocol categories, and ten-minute precedent (R06); learned normal
  device behaviour (R11--R17); adaptive learning and contamination concerns
  (R18--R20). Attribute these before describing the implementation.
\item \textbf{Project choices requiring local justification:} the joint minimum
  of standardised DM, total-volume, and ARP scores; the selected quantile;
  the exact confidence margins, buffer lengths, refit schedules, movement
  bounds, fixed-reference decisions, and persistence/rollback design.
  Related literature motivates the problems; none of the reviewed papers
  validates this exact combination or proves it novel.
\item \textbf{Overdispersion:} compare conditional count variation with the
  multinomial expectation at the same total count. ``Variance exceeds the
  mean'' is not a universal definition for multinomial counts.
  Define the project's inverse concentration $\psi=1/\alpha_0$ explicitly
  rather than assuming all sources use the same dispersion parameter.
\item \textbf{Modelling versus detection:} the DM is conditional on the
  observed packet total. A good category-profile fit does not establish good
  attack detection or make raw likelihood an attack probability. Separate
  count margins are a project design choice. Residual explanations show
  statistical deviation, not causal attack attribution.
\item \textbf{Unlabelled learning versus unlabelled research:} attack labels
  are absent from fitting and runtime update admission. They are present
  in offline evaluation and have informed development comparisons. A later
  evaluation partition with prior label exposure is not an untouched final test.
\item \textbf{Measured outcomes:} precision, recall, alert delay, adaptation
  improvements, numerical error, and resource usage require this project's
  own versioned evidence. None can be borrowed from these publications.
\end{itemize}

\subsection{Existing resources deliberately left out of the main shortlist}

\path{resouces/Proposte_di_Ricerca_basate_su_Dirichlet_Multinomial.pdf} is
the local project brief, useful for scope and provenance rather than external
scientific evidence. Its neural DirichletSoftmax proposal is a different
project, not part of the implemented detector.

The locally supplied survey \emph{Optimal Countermeasures Selection Against
Cyber Attacks: A Comprehensive Survey on Reaction Frameworks} concerns response
and countermeasure selection. It could support a future-work paragraph on
automated reaction, but LM-IDNet emits alerts rather than selecting or enforcing
countermeasures. It was not used merely to increase the reference count.

\subsection{Using the files}

This review compiles independently of the thesis:
\begin{verbatim}
./scripts/latex.sh docs/thesis_reference_review.tex
\end{verbatim}
The PDF is written to \path{out/thesis_reference_review.pdf}. The build uses
the repository's existing TeX Live/latexmk service and Biber for the companion
\path{docs/thesis_reference_review.bib}.

The companion file contains all 30 scholarly records and D1--D2 with stable
citation keys. Copy selected entries into \path{thesis/references.bib} as the
corresponding passages are written. For example, cite
\verb|\cite{minka2000,languasco2023}| when distinguishing estimation from
likelihood evaluation, and \verb|\cite{hamza2019,unswattackdata}| when
introducing the attack traces. Read the cited passage before final submission,
especially the three limited-access entries and version-sensitive R07.

\end{document}
